\documentclass[aspectratio=169,usenames,dvipsnames]{beamer}
\beamertemplatenavigationsymbolsempty
% \documentclass{beamer}

% UNCOMMENT FOR HANDOUTS
% \usepackage{pgfpages}
% \pgfpagesuselayout{2 on 1}[landscape,letterpaper,border shrink=5mm]
% \setbeameroption{show notes}

\mode<presentation>
\usetheme{Boadilla}
\usecolortheme{spruce}
\usepackage{../LizStyle,ulem}
\usepackage{color}
\usepackage[color,matrix,arrow]{xy}
% \input xy
% \xyoption{all}
\usepackage{tikz}
\usepackage{tikz-cd}
\tikzset{commutative diagrams/diagrams={ampersand replacement=\&}}

\AtBeginSection{\frame{\sectionpage}}


% ==========Command for putting comments and removing them=======
% Visible:
\newcommand{\InstructorNotes}[1]{\textit{\color{orange}#1}}
\newcommand{\InstructorList}[1]{
\begin{itemize}
	\color{orange}
	#1
\end{itemize}
}
% Removed:
% \renewcommand{\InstructorNotes}[1]{}
% \renewcommand{\InstructorList}[1]{}


\newcommand{\bvfill}{\vskip0pt plus 1filll}

\usepackage{multimedia}
\usepackage{graphbox}

% \usepackage{algpseudocode}


\newcommand{\Reeb}{\textbf{Reeb}}
\newcommand{\Rtopc}{\textbf{Rtopc}}
% \newcommand{\RR}{\mathcal{R}}

\newcommand{\Set}{\textbf{Set}}
\newcommand{\Vect}{\textbf{Vect}}
\newcommand{\Top}{\textbf{Top}}
\newcommand{\Int}{\textbf{Int}}


\graphicspath{ {Figures/} {../Figures/}{../../Figures/} }



\begin{document}
%\pdfoptionpdfminorversion 5
%\logo{\includegraphics[height=1.5cm]{dukeshield.pdf}}
\title[]{Ch 6.3: PCR}
\subtitle{Lecture 20 - CMSE 381}
\date{Fri, Oct 18, 2024}



\institute[MSU-CMSE]{Michigan State University\\
			::\\
          Dept of Computational Mathematics, Science \& Engineering}
\author[Dr.~Munch]{Prof.~Elizabeth Munch}

\begin{frame}
 \titlepage


\end{frame}
%--------------------------
\begin{frame}[c]{Announcements}

\begin{columns}
\column{.45\textwidth}
\centering

	\textbf{Last time:}
	\begin{itemize}
		\item PCA
	\end{itemize}

	\vspace{.2in}

	\textbf{This lecture:}
	\begin{itemize}
	\item PCR
\end{itemize}


\textbf{Announcements:}
	\begin{itemize}
		\item Exam \#2 on Friday!
		\begin{itemize}
			\item Bring 8.5x11 sheet of paper 
			\item Handwritten both sides 
			\item Anything you want on it, but must be your work 
			\item You will turn it in
			\item Non-internet calculator if you want it
		\end{itemize} 

	\end{itemize}

\column{.45\textwidth}
\centering
\includegraphics[width = \textwidth]{Day01-Schedule-Test2.png}
\end{columns}

\end{frame}
%--- Next Frame ---%

%=================
\section{Last time}
%=================


%--- Next Frame ---%
\begin{frame}[t]{Shrinkage}
	Find $\beta$ to minimize
	\begin{equation*}
		RSS = \sum_{i=1}^n \left(y_i - \beta_0 - \sum_{j=1}^p \beta_j x_{ij}\right)^2
	\end{equation*}
	subject to:

\begin{columns}
\column{.3\textwidth}
\centering
\textbf{Least Squares:}\\
\footnotesize
No constraints

\column{.3\textwidth}
\centering
\textbf{Ridge:}
\footnotesize
\begin{equation*}
	 \sum_{j=1}^p \beta_j^2 \leq s
\end{equation*}

\column{.3\textwidth}
\centering

\textbf{The Lasso:}
\footnotesize
\begin{equation*}
	 \sum_{j=1}^p |\beta_j| \leq s
\end{equation*}

\end{columns}


\begin{columns}
\column{.4\textwidth}
\centering
\includegraphics[width = \textwidth]{Fig6-4}
\column{.4\textwidth}
\centering

\includegraphics[width = \textwidth]{Fig6-6}
\end{columns}


\InstructorNotes{Also, find best choice of $\lambda$ or $s$ using CV}

\end{frame}
%--- Next Frame ---%


\begin{frame}[t]{Linear transformation of predictors }

\begin{columns}
\column{.45\textwidth}
\centering
\textbf{Original Predictors:}
\begin{equation*}
	X_1,\cdots,X_p
\end{equation*}
\textbf{New Predictors:}
\begin{equation*}
	Z_1,\cdots,Z_M
\end{equation*}
\begin{equation*}
	Z_m = \sum_{j=1}^p \phi_{jm}X_j
\end{equation*}
\column{.45\textwidth}
\centering
\textbf{The goal:}
\begin{itemize}
\item Find good $\phi$'s (PCA)
\item Fit regression model on $Z_i$'s using least squares (PLS)
\begin{equation*}
	y_i = \theta_0 + \sum_{m=1}^M \theta_m z_{im} + \e_i
\end{equation*}
\item Hope that lower dimensions means less overfitting
% \item Remember that interpretation not the same as shrinkage/subset selection of variables
\end{itemize}
\end{columns}

\end{frame}
%--- Next Frame ---%
% \begin{frame}[c]{PCA - First PC}

% \begin{columns}
% \column{.4\textwidth}
% \centering
% \includegraphics[width = \textwidth, align = c]{PCA-ex-A}
% Projection has high variance
% \column{.4\textwidth}
% \centering
% \includegraphics[width = \textwidth, align = c]{PCA-ex-B}
% Projection has low(er) variance

% \end{columns}

% \end{frame}
%--- Next Frame ---%
\begin{frame}[t]{PCA - First PC}
	\begin{columns}
	\column{.24\textwidth}
	\includegraphics[align = c, width = \textwidth]{DimReduction_Projection_Example2_78.png} 
	\includegraphics[align = c, width = \textwidth]{DimReduction_Projection_Example2_hist78.png} 
	\column{.24\textwidth}
	\includegraphics[align = c, width = \textwidth]{DimReduction_Projection_Example2_157.png} 
	\includegraphics[align = c, width = \textwidth]{DimReduction_Projection_Example2_hist157.png} 
	
	\column{.24\textwidth}
	\includegraphics[align = c, width = \textwidth]{DimReduction_Projection_Example2_0.png} 
	\includegraphics[align = c, width = \textwidth]{DimReduction_Projection_Example2_hist0.png} 

	\column{.24\textwidth}
	\includegraphics[align = c, width = \textwidth]{DimReduction_Projection_Example2_235.png} 
	\includegraphics[align = c, width = \textwidth]{DimReduction_Projection_Example2_hist235.png} 


	\end{columns}

	\InstructorNotes{higher variance of projection on the further right pic}
	\end{frame}
	%--- Next Frame ---%

\begin{frame}[c]{Projection onto first PC}

	\centering
	\includegraphics[width =.6 \textwidth, align = c]{Fig6-15a}

	\begin{equation*}
		Z_1 = 0.839 \cdot (\texttt{pop} - \overline{\texttt{pop}}) + 0.544 \cdot (\texttt{ad} - \overline{\texttt{ad}})
	\end{equation*}
	\InstructorNotes{$\phi_{11} = 0.839$ and $\phi_{21} = 0.544$ are the principal component loadings. Blue dot is $(\overline{pop},\overline{ad})$ }


\end{frame}
%--- Next Frame ---%




\begin{frame}[t]{Drawing points in PC space}

\centering
\includegraphics[width = .8\textwidth, align = c]{Fig6-15}
\begin{columns}
\column{.45\textwidth}
\centering
\column{.45\textwidth}
\centering

\end{columns}

\end{frame}
%--- Next Frame ---%

%==========
\section{Principal Components Regression }
%==========

\begin{frame}[t]{So you've found your PCA coefficients}

\begin{columns}
\column{.45\textwidth}
\centering
Now what?

\InstructorList{
\item Do linear regression on $Z_1,\cdots,Z_M$
\item Book calls this "principal components regression "
\item Linear model is $y = \theta_0 + \theta_1 Z_1 + \cdots + \theta_M Z_M$
}
\column{.45\textwidth}
\centering
What are we assuming?

\InstructorList{
\item The directions in which $X_1,\cdots, X_p$ have the most variation are the directions that are associated with $Y$
\item If assumption holds, then fitting on $Z_i$ better than fitting on $X_i$ since fewer variables lessens chance of overfitting
}

\end{columns}

\end{frame}
%--- Next Frame ---%
\begin{frame}[t]{Interpretation of PCR coefficients }
	\begin{columns}
	\column{.35\textwidth}
	\centering
	\textbf{Original Predictors:}
	\begin{equation*}
		X_1,\cdots,X_p
	\end{equation*}
	\textbf{New Predictors:}
	\begin{equation*}
		Z_1,\cdots,Z_M
	\end{equation*}
	\begin{equation*}
		Z_m = \sum_{j=1}^p \phi_{jm}X_j
	\end{equation*}
	\textbf{Learned model:}
	\begin{equation*}
		y = \theta_0 + \theta_1 Z_1 + \cdots + \theta_M Z_M
	\end{equation*}
	
	
	
	\column{.55\textwidth}
	\centering
	\footnotesize
	\InstructorList{
		\item Do example showing that coefficients of $Z_i$ can't be used to say anything about $X_i$'s 
		\item $p=3, M=2$
		\item 
		$\phi = \begin{pmatrix}
		1 & -1  \\
		3 & 1/2 \\
		0 & 1/2
		\end{pmatrix}$
		\item  Learned model $ y = 2 Z_1 - Z_2$
		\item $y = 2(X_1+3X_2 + 0X_3) - (-X_1 + \tfrac{1}{2} X_2 +  \tfrac{1}{2} X_3  )$
		\item $= 3X_1 + 2.5 X_2 - 0.5 X_3 $
		\item Conclusion: can't directly use coefficients learned to interpret original inputs! 
	  \item Each PC is a linear combination of the originals, so can go backwards to figure out whaat the coefficients would have been, see left figure 
	}
	\end{columns}
	\end{frame}
%--- Next Frame ---%
%--- Next Frame ---%
\begin{frame}[c]{Picking $M$}

	\begin{columns}
	\column{.75\textwidth}
	\centering
	\includegraphics[width = \textwidth]{Fig6-20}
	% \column{.45\textwidth}
	% \centering
	
	\end{columns}
	\InstructorList{
	\item Credit data set
	\item Major drop in MSE at 10 compoentns, so barely a savings with the original 11 variables
	}
	
	\end{frame}
	%--- Next Frame ---%	
{
\setbeamercolor{background canvas}{bg=Apricot!25}
\begin{frame}[t]{Do PCR with hitters data}
\begin{columns}
\column{.45\textwidth}
\centering
\column{.45\textwidth}
\centering
\end{columns}
\end{frame}
}
%--- Next Frame ---%

\begin{frame}[c]{Doing better}
Example with simulated data: $n=50$ observations of $p=45$ predictors\\
\begin{columns}
\column{.45\textwidth}
\centering
\InstructorList{
\item Right side is usual least squares estimate
\item Typical U shape for Test MSE
\item PCR with good choice of number of components improvement over plain least squares
}
\column{.45\textwidth}
\centering
$Y$ is a function of \textbf{2 predictors}
\includegraphics[width = \textwidth, align = c]{Fig6-18b}
\end{columns}

\end{frame}
%--- Next Frame ---%
\begin{frame}[c]{Doing better}
Example with simulated data: $n=50$ observations of $p=45$ predictors\\
\begin{columns}
\column{.45\textwidth}
\centering
$Y$ is a function of \textbf{all predictors}
\includegraphics[width = \textwidth, align = c]{Fig6-18a}
\column{.45\textwidth}
\centering
\InstructorList{
\item Right side is usual least squares estimate
\item Even better improvement using PCR than previous
\item note the chnage in $y$-axis values
}
\end{columns}

\end{frame}
%--- Next Frame ---%

\begin{frame}[c]{Comparison to results on shrinkage}

\begin{columns}
\column{.45\textwidth}
\centering
$Y$ is a function of all predictors
\includegraphics[width = .5\textwidth, align = c]{Fig6-18a}
\includegraphics[width = \textwidth, align = c]{Fig6-8}
\column{.45\textwidth}
\centering
$Y$ is a function of 2 predictors
\includegraphics[width = .5\textwidth, align = c]{Fig6-18b}
\includegraphics[width = \textwidth, align = c]{Fig6-9}

\end{columns}
\InstructorNotes{Shrinkage does better}
\end{frame}


%--- Next Frame ---%
\begin{frame}[c]{Properties of PCR}

	\begin{columns}
	\column{.45\textwidth}
	\centering
	\InstructorList{
	\item Note better in previous example because lots of PCs needed
	\item It does better when info contained in the first few PCs
	\item Number of PCs to use can be picked with CV
	\item Standardization needed to make sure big data values don't skew the results
	}
	\column{.45\textwidth}
	\centering
	\InstructorList{
	\item Not a feature selection method because we don't get a subset of features, we get a collection of linear combos of features
	}
	\end{columns}
	
	\end{frame}
	%--- Next Frame ---%
% \begin{frame}[c]{Principal Components Regression (PCR)}

% \begin{columns}
% \column{.4\textwidth}
% \centering
% \begin{itemize}
% 	\item Take new features $Z_i$
% 	\item Run regression
% 	\item Maybe do CV for a bunch of choices of $M$ (where $M = $ number of $Z_i$ features used) to pick a best $M$
% \end{itemize}
% \column{.5\textwidth}
% \centering
% \InstructorList{
%   \item Note that you can't directly use coefficients learned to interpret original inputs 
%   \item Each PC is a linear combination of the originals, so can go backwards to figure out whaat the coefficients would have been, see left figure 
% }

% \end{columns}

% \end{frame}

% %--- Next Frame ---%
% \begin{frame}[c]{Bias variance trade-off}
% Example with simulated data: $n=50$ observations of $p=45$ predictors\\
% $Y$ is a function of \textbf{all predictors}
% \begin{columns}
% \column{.45\textwidth}
% \centering
% \includegraphics[width = \textwidth, align = c]{Fig6-18a}
% \column{.45\textwidth}
% \centering
% \InstructorList{
% \item Right side is usual least squares estimate
% \item Even better improvement using PCR than previous
% \item note the chnage in $y$-axis values
% }
% \end{columns}

% \end{frame}
% %--- Next Frame ---%
% \begin{frame}[c]{Bias variance trade-off}
% Example with simulated data: $n=50$ observations of $p=45$ predictors\\
% $Y$ is a function of \textbf{2 predictors}
% \begin{columns}
% \column{.45\textwidth}
% \centering
% \InstructorList{
% \item Right side is usual least squares estimate
% \item Typical U shape for Test MSE
% \item PCR with good choice of number of components improvement over plain least squares
% }
% \column{.45\textwidth}
% \centering
% \includegraphics[width = \textwidth, align = c]{Fig6-18b}
% \end{columns}

% \end{frame}
%--- Next Frame ---%

% \begin{frame}[c]{Comparison to results on shrinkage}

% \begin{columns}
% \column{.45\textwidth}
% \centering
% $Y$ is a function of all predictors
% \includegraphics[width = .5\textwidth, align = c]{Fig6-18a}
% \includegraphics[width = \textwidth, align = c]{Fig6-8}
% \column{.45\textwidth}
% \centering
% $Y$ is a function of 2 predictors
% \includegraphics[width = .5\textwidth, align = c]{Fig6-18b}
% \includegraphics[width = \textwidth, align = c]{Fig6-9}

% \end{columns}
% \InstructorNotes{Shrinkage does better}
% \end{frame}
%--- Next Frame ---%
% \begin{frame}[c]{Picking $M$}

% \begin{columns}
% \column{.75\textwidth}
% \centering
% \includegraphics[width = \textwidth]{Fig6-20}
% % \column{.45\textwidth}
% % \centering

% \end{columns}
% \InstructorList{
% \item Credit data set
% \item Major drop in MSE at 10 compoentns, so barely a savings with the original 11 variables
% }

% \end{frame}





%--- Next Frame ---%
	

% {
% \setbeamercolor{background canvas}{bg=Apricot!25}
% \begin{frame}[t]{Figuring out the original model from the PC model }
% \begin{columns}
% \column{.45\textwidth}
% \centering
% \begin{itemize}
% 	\item 
% We have three input variables $X_1, X_2, X_3$
% \item 
% We have two PCs, 
% \begin{align*}
% 	Z_1 &= 0.266 \cdot X_1 -0.077 \cdot X_2 + 0.961 \cdot X_3\\ 
% 	Z_2 &= 0.968 \cdot X_1 +0.136 \cdot X_2 + -0.254 \cdot X_3 
% \end{align*}
% \item Using linear regression, we learn the model 
% \begin{equation*}
% 	Y = -3 + 2Z_1 -4 Z_2.
% \end{equation*}
% \item What are the coefficients for the model in terms of the $X_i$'s? 
% \end{itemize}
% \column{.45\textwidth}
% \centering
% \InstructorNotes{
% 	\tiny
%   \begin{align*}
% 	Y &= -3 + 2Z_1 -4 Z_2\\
% 	&= -3 + 2(0.266 \cdot X_1 -0.077 \cdot X_2 + 0.961 \cdot X_3)\\
% 	& \qquad -4(0.968 \cdot X_1 +0.136 \cdot X_2 + -0.254 \cdot X_3 )\\
% 	&= -3 - 3.34 X_1 - 0.698 X_2 + 2.938 X_3
%   \end{align*}
% }

% \end{columns}
% \end{frame}
% }
% %--- Next Frame ---%
% \begin{frame}[t]{Example on Credit dataset}
% \begin{columns}
% \column{.6\textwidth}
% \includegraphics[width = \textwidth, align = c]{Fig6-20}
% \centering
% % \column{.45\textwidth}
% \centering
% \end{columns}

% \InstructorList{
%   \item Not a subset selection method because all $X_i$'s still seen 
%   \item Book says you can view ridge regression as a continuous version of PCR
%   \item In this example, would likely pick $M=10$ components.  Since had 11 originally, not much reduction
% }

% \end{frame}
%--- Next Frame ---%




% \begin{frame}[t]{First direction $Z_1$ for Partial Least Squares (PLS)}

% 	\begin{columns}
% 	\column{.45\textwidth}
% 	\centering
% 	\begin{itemize}
% 		\item Set $\phi_{j1}$ equal to the coefficient from simple linear regression of $Y$ onto $X_j$
% 		\InstructorNotes{ \footnotesize
% 	 This means we do $p$ linear regressions, not that we take the multivariable linear regression coefficients
% 		}
% 		\item The first direction is
% 		\begin{equation*}
% 			Z_1 = \sum_{j=1}^p \phi_{j1}X_j
% 		\end{equation*}
% 	\end{itemize}
	
% 	\footnotesize
% 	\InstructorList{
% 	\item
% 	PLS places highest weight on variables most strongly related to the response
% 	\item In example, PLS chooses direction less change in ad direction than pop relative to PCA
% 	\item Implies pop more highly correlated with response than ad
% 	}
	
% 	\column{.45\textwidth}
% 	\centering
% 	\includegraphics[width = \textwidth]{Fig6-21}
% 	Ex.~Prediction of $Y=$Sales (not shown) on $X_1=$Population and $X_2=$Ad Spending
% 	\begin{itemize}
% 		\item Solid green: First PLS direction
% 		\item Dashed: First PC direction
% 	\end{itemize}
% 	\end{columns}
	
% 	\end{frame}
% 	%--- Next Frame ---%
	
% 	\begin{frame}[c]{Second (and more) PLS directions}
	
% 	\begin{columns}
% 	\column{.45\textwidth}
% 	\centering
% 	\begin{itemize}
% 		\item Regress each variable on $Z_1$ and take residuals
% 		\InstructorNotes{This is the remaining info not explained by first PLS direction}
% 		\item Compute $Z_2$ using \textit{orthogonalized data} same as for $Z_1$
% 	\end{itemize}
% 	\InstructorList{
% 	\item Number $M$ can be picked using CV after standardizing predictors
% 	\item Mixed bag in terms of results..... in practice often performs no better than ridge or PCR
% 	}
% 	\column{.45\textwidth}
% 	\centering
	
% 	\end{columns}
	
% 	\end{frame}
% 	%--- Next Frame ---%
	
	
	
	\begin{frame}[c]{TL;DR}
	
	\begin{columns}
	\column{.45\textwidth}
	\centering
	\textbf{PCR}
	\begin{itemize}
		\item Unsupervised dimensionality reduction $+$ linear regression
		\item Choose component $Z_1$ in the direction of most variance using only $X_i$'s information
		\item Choose $Z_2$ and beyond by the same method after ``getting rid'' of info in the directions already explained
	\end{itemize}
	\includegraphics[width = .5\textwidth, align = c]{Fig6-15a}
	
	\column{.45\textwidth}
	\centering
	% \textbf{PLS}
	% \begin{itemize}
	% 	\item Supervised dimensionality reduction
	% 	\item Choose component $Z_1$ by using simple regression coefficients of each $X_i$ onto $Y$
	% 	\item Choose $Z_2$ and beyond by the same method after ``getting rid'' of info in the directions already explained
	% 	\item Not a particular benefit, so usually default to PCA unless you have a good reason for this
	% \end{itemize}
	% \includegraphics[width = .5\textwidth]{Fig6-21}
	
	\end{columns}
	
	\end{frame}
	%--- Next Frame ---%
	

\begin{frame}[t]{Next time}



\includegraphics[align = c, width = .5\textwidth]{Day01-Schedule-Test2}

\end{frame}
%--- Next Frame ---%



\end{document}
